\documentclass[11pt]{article}
\usepackage[margin=25mm]{geometry}
\usepackage{fontspec}\setmainfont{Latin Modern Roman}
\usepackage{amsmath,amssymb,mathtools}
\usepackage{unicode-math}\setmathfont{Latin Modern Math}
\usepackage{xurl,hyperref}
\hypersetup{colorlinks=true,urlcolor=blue}
\setcounter{secnumdepth}{0}\setlength{\parindent}{0pt}\setlength{\parskip}{5pt}\setlength{\emergencystretch}{3em}
\newcommand{\sref}[2]{\hyperlink{source:#1:#2}{[S#2]}}
\newcommand{\cataloguescope}[1]{\paragraph{Recorded target.}#1}
\begin{document}
\section{Kourovka Notebook Problem 21.99}\label{statement}

\subsection{Definitions and mathematical statement}

\paragraph{Shared notation.}$\mathbb N=\{1,2,\ldots\}$, and $\mathbb Z,\mathbb Q,\mathbb R,\mathbb C$ denote the integers, rationals, reals and complex numbers. Any other conventions are specified locally.


Let a finite group $G\leq\operatorname{Sym}(\Omega)$ act transitively on a finite set $\Omega$. For $g\in G$, let $\operatorname{Fix}_\Omega(g)=\{\omega\in\Omega:g\omega=\omega\}$. A derangement has no fixed points, but the requested conclusion also permits elements with at least two fixed points.
\paragraph{Statement recorded in the supplied source.}
For every such action and every distinct $\alpha,\beta\in\Omega$, is there a $g\in G$ such that
\[
 g\alpha=\beta\qquad\text{and}\qquad|\operatorname{Fix}_\Omega(g)|\ne1?
\]
No primitivity assumption or requirement that $g$ be a derangement is imposed. \sref{statement}{2}
\subsection{Short English statement}
Can every ordered pair of distinct points in a finite transitive action be joined by an element having either zero or at least two fixed points?
\subsection{Sources}
\begin{description}
\item[\hypertarget{source:statement:1}{S1}] ULAM.AI and the UnsolvedMath Contributors, \emph{UnsolvedMath: A Curated Collection of Open Mathematics Problems}, catalogue entry. CC BY 4.0. \url{https://huggingface.co/datasets/ulamai/UnsolvedMath}.
\item[\hypertarget{source:statement:2}{S2}] P. Müller, Problem 21.99, in E. I. Khukhro and V. D. Mazurov (eds.), \emph{The Kourovka Notebook: Unsolved Problems in Group Theory}, 21st edition, arXiv:1401.0300v48 (6 October 2026). This version records a refutation; the displayed question is the historically registered target. \url{https://arxiv.org/abs/1401.0300v48}.
\end{description}
\label{end:statement}
\end{document}
